
% Chapter 2 File

\chapter{Literature Review}
\label{chapter2}
\thispagestyle{empty}

Many authors previously covered customer behaviour on online marketplaces, comparing it with the offline domain (e.g. Brynjolfsson and Smith, 2000; Brynjolfsson et al., 2003). Recently, literature also focused on user behaviour on mobile app stores: for example, Burghers et al. (2016, p. 330) argued how apps are generally low-priced or free and this causes low effort and regret when they are unsatisfied and decide to uninstall them. This means that the perceived risk in the buying decision process is low and leads, in many cases, to low involvement. In addition, some consumers appear to be unconscious and they do not seem to search for relevant information that could reduce the information asymmetries in the app market (Buck et al., 2014, p. 12). There is evidence that decision-making when dealing with app downloading often relies on heuristics (e.g. ``take the first"), which allows users to quickly navigate across offerings, but may also cause customers to overlook some important information (Dogruel et al., 2015, p. 139).

In this context, it can be assumed that some elements in the presentation page of an app on a store may be effective in convincing a user to download it. In particular, three main references and their influence on user's intention to download will be investigated below: app reputation, its popularity (and the consequent peer influence) and the developer's brand.

Product reputation impact on purchase decisions online has already been covered by many authors. Chevalier and Mayzlin (2006, pp. 353-354), for example, proved that an improvement in online reviews is positively correlated with an increase in sales and the impact of ratings isn't linear, with one-star reviews being more influential than five-star reviews. Moe and Trusov (2011, pp. 455-456) testified how ratings can improve sales, but customer rating behaviour is significantly socially influenced by previously posted ratings. Previous literature investigated eWOM (electronic word-of-mouth) valence, usually reflected by the average rating score, as well as eWOM volume, which is usually measured by the number of reviews. Out of the papers that covered eWOM valence on mobile app stores, seven of them were able to find a positive impact on app success (Böhm and Schreiber, 2014, p. 13; Burgers et al., 2016, p. 340; Chen et al., 2022, pp. 6-7; Finkelstein et al., 2017, p. 38; Gokgoz et al., 2021, pp. 434-435; Jung et al., 2012, p. 936; Lee and Raghu, 2014, p. 161), one found a negative correlation (Liang et al., 2015, pp. 250-251) and four found no significant impact (Roma and Ragaglia, 2016, pp. 184-185; Sällberg et al., 2022, pp. 7-10; Timmerman and Shepherd, 2016, pp. 14-15; Wang et al., 2015, p. 5). On the other side, four authors discovered a positive relation between app success and the number of reviews (Burgers et al., 2016, p. 340; Lee and Raghu, 2014, pp. 149-151; Sällberg et al., 2022, p. 13; Wang et al., 2015, pp. 7-10), while just one found out it was negative (Liang et al., 2015, pp. 250-251). The majority, in both cases, used the number of downloads in a selected time period as a performance metric, while some others preferred app survival in the top-apps lists. Overall, reputation in the mobile ecosystem has already been studied enough to take some preliminary conclusions and the majority of past literature contributions agree that reputation signals quality and it has a generally positive effect on sales.

The popularity of an app is reflected in multiple elements of an app store: firstly, the number of downloads is usually shown in the app page, then the top-list presence and app sales rank also signal it is popular among users (Sällberg et al., 2022, p. 3). Popularity translates into more interest among potential customers because it reflects the previous interest of other consumers and this peer influence has been shown to be one of the primary drivers of purchase decisions in online shopping (Liu et al., 2014, p. 330). Salganik et al. (2006, p. 855) have already shown how popularity information and social influence affect other users' choices on music platforms, making popular songs even more popular. Also in the software industry, there is evidence that app popularity, reflected, for example, in the presence in top-apps ranking, has a positive impact on app success (e.g. Finkelstein et al., 2017, p. 2679). In addition, his effect may be even stronger in the case of software that generates strong network effects, since its utility increases with popularity and users are expected to consider that. Finally, the consequences of reputation aren't independent of social influence: interestingly, less popular products in the market are more significantly influenced by online reviews (Liu et al., 2014, p. 331).

The third and last considered anchor is the brand of the company developing the app. Even if online marketplaces lower prices and reduce information asymmetries, brands remain sources of competitive advantage (e.g. Brynjolfsson and Smith, 2000, p. 580). In the mobile ecosystem, Roma and Ragaglia (2016, p. 184), for example, have already shown how a company's fame positively moderates app success. On mobile app stores, the brand may be reflected in multiple elements of the presentation page: apart from the developer or publisher name, which is usually present, the company's name is sometimes included also in the app's name and in the icon. While, if taken singularly, each one of these elements has a relative importance on purchase decision which is low, at least if compared to average review rating and number of reviews (Böhm and Schreiber, 2014, p. 12), they could combine and consistently influence users. Furthermore, app icons are graphic elements and they are easy to understand for users (Böhm and Schreiber, 2014, p. 8) and Pol (2015, p. 51) showed how consumers prefer apps with a brand logo, well-known brands also gain an additional advantage and the brand attitude has a positive effect on the perceived app quality.

